%! Logarithmic Function Context and Data Modeling — Unit 2, Lesson 2.14 — Guided Notes (Answer Key)
\documentclass[10pt]{article}
\usepackage{precalculus-article}
\usepackage{precalculus-key}   % pulls in -boxes; do NOT also load -boxes
\newcommand{\TallMath}[1]{$\displaystyle #1\rule[-1.4em]{0pt}{3.2em}$}

\begin{document}

\pageheader{Unit 2, Lesson 2.14}{Guided Notes --- Answer Key}
\namedateperiod

\vspace{0.1in}

\begin{objectivebox}
\textbf{LO 2.14.A:} Construct a logarithmic function model.

By the end of this lesson I will be able to:
\begin{itemize}[leftmargin=1.5em,itemsep=3pt,topsep=3pt]
  \item Recognize when a context or data set calls for a logarithmic model.
  \item Build a logarithmic model from two input-output pairs using a system of equations.
  \item Apply transformations of $f(x) = a\log_b x$ to fit a model to a context.
  \item Use a logarithmic model to predict values of the dependent variable.
\end{itemize}
\end{objectivebox}

\vspace{0.08in}

\begin{vocabbox}
\termblanklong{logarithmic model}
\ansline{A function of the form $f(x)=a\log_b x + k$ used when inputs change proportionally
over equal-length output intervals.}

\termblanklong{proportion}
\ansline{The constant factor by which the input is multiplied as the output increases by one
unit; equals the base $b$ of the logarithmic model.}

\termblanklong{real zero}
\ansline{The input value where the model output equals zero; for $f(x)=\log_b x$ the real
zero is $x=1$.}

\termblanklong{natural logarithm model}
\ansline{A logarithmic model using base $e$: $f(x)=a\ln(x)+k$.}

\termblanklong{logarithmic regression}
\ansline{A technology-assisted technique that fits a logarithmic curve to a data set.}

\termblanklong{prediction}
\ansline{Using a model to estimate an output for a given input, or to find the input that
produces a given output.}
\end{vocabbox}

\vspace{0.08in}

\begin{hookbox}
\textbf{Consider:} A scientist observes that a bacterial population doubles in the 1st hour,
doubles again after 2 more hours, then again after 4 more hours.

\vspace{4pt}
How much time does each additional doubling require?
\ans{1 hr}, \ans{2 hr}, \ans{4 hr}

\vspace{4pt}
As the population grows, the time per doubling is \ans{increasing}.

\vspace{4pt}
Why might a logarithmic function --- not an exponential one --- model this situation?
\ansline{Equal output increments (one more doubling) correspond to proportionally growing
input values (the time intervals multiply by 2), which is the defining pattern of a
logarithmic model.}
\end{hookbox}

\vspace{0.08in}

\begin{notesbox}{When to Use a Logarithmic Model (EK 2.14.A.1)}
A logarithmic model is appropriate when input values change \ans{proportionally (multiply
by a constant)} over equal-length output-value intervals.

\vspace{6pt}
\textbf{Contrast:}
\begin{itemize}[leftmargin=1.5em,itemsep=4pt,topsep=2pt]
  \item \textbf{Exponential:} output changes \ans{proportionally} over equal input intervals.
  \item \textbf{Logarithmic:} input changes \ans{proportionally} over equal output intervals.
\end{itemize}

\vspace{8pt}
\textbf{Example.} Fill in the table and answer the questions.

\vspace{4pt}
\begin{center}
\renewcommand{\arraystretch}{1.8}
\begin{tabular}{|c|c|}
\hline
\rowcolor{sky}
$x$ (area, hectares) & $y$ (number of species) \\
\hline
$2$  & $1$ \\
\hline
$4$  & $2$ \\
\hline
$8$  & $3$ \\
\hline
$16$ & $4$ \\
\hline
\end{tabular}
\end{center}

\vspace{6pt}
As $y$ increases by 1, $x$ \ans{doubles (multiplies by 2)}. \quad Proportion: \ans{2}

\vspace{4pt}
This pattern suggests the model: $y = \log_{\ans{2}}(x)$
\end{notesbox}

\vspace{0.08in}

\begin{notesbox}{Building a Model from Two Input-Output Pairs (EK 2.14.A.2)}
\textbf{General form:} $f(x) = a\log_b x + k$

\vspace{6pt}
\textbf{Worked Example.} Given: $f(1) = 0$ and $f(e) = 3$. Find the model.

\vspace{4pt}
\textit{Step 1.} Substitute $(1,\,0)$: \quad $0 = a\ln(1) + k = a \cdot \ans{0} + k$
\quad so $k = $ \ans{0}

\vspace{4pt}
\textit{Step 2.} Substitute $(e,\,3)$: \quad $3 = a\ln(e) + 0 = a \cdot \ans{1}$
\quad so $a = $ \ans{3}

\vspace{4pt}
\textit{Model:} $f(x) = $ \ans{$3\ln(x)$}

\vspace{8pt}
\textbf{Practice.} Given: $f(1) = 2$ and $f(9) = 4$. Use base $b = 3$.

\vspace{4pt}
\textit{Step 1.} From $f(1) = 2$: since $\log_3(1) = $ \ans{0}, we get $k = $ \ans{2}

\vspace{4pt}
\textit{Step 2.} From $f(9) = 4$: $4 = a\log_3(9) + k = a \cdot \ans{2} + \ans{2}$
\quad so $a = $ \ans{1}

\vspace{4pt}
\textit{Model:} $f(x) = $ \ans{$\log_3(x) + 2$}
\end{notesbox}

\vspace{0.08in}

\begin{notesbox}{Transformations of a Logarithmic Model (EK 2.14.A.3)}
General form with transformations: $f(x) = a\log_b(x - h) + k$

\vspace{4pt}
\begin{itemize}[leftmargin=1.5em,itemsep=4pt,topsep=2pt]
  \item $a$: vertical \ans{dilation} (stretch if $|a|>1$, compression if $|a|<1$)
  \item $h$: horizontal shift \ans{$h$} units to the \ans{right}
  \item $k$: vertical shift \ans{$k$} units \ans{up}
\end{itemize}

\vspace{6pt}
\textbf{Context example.} A sound-level function has output 0 dB at intensity $x = 1$
and increases by 10 dB for every tenfold increase in intensity.

\vspace{4pt}
The output is 0 when $x = 1$, so \ans{$k$} (which parameter?) equals \ans{0}.

\vspace{4pt}
A tenfold increase in $x$ adds 10 to the output, so $a = $ \ans{10} and $b = $ \ans{10}.

\vspace{4pt}
Model: $S(x) = $ \ans{$10\log_{10}(x)$}
\end{notesbox}

\vspace{0.08in}

\begin{notesbox}{Predicting with a Logarithmic Model (EK 2.14.A.6)}
Given $f(x) = 2\log_3 x$:

\vspace{6pt}
\textbf{Predict $f(81)$:} \quad $f(81) = 2\log_3 81 = 2 \cdot \ans{4} = $ \ans{8}

\vspace{6pt}
\textbf{Solve for $x$ when $f(x) = 5$:}

\vspace{4pt}
$2\log_3 x = 5 \;\Rightarrow\; \log_3 x = \ans{5/2} \;\Rightarrow\; x = 3^{\ans{5/2}} = $ \ans{$9\sqrt{3}$}

\vspace{8pt}
\textbf{Natural log models (EK 2.14.A.5).} Recall $\ln x = \log_e x$.

\vspace{4pt}
Useful identities: \quad $\ln(e^n) = $ \ans{$n$} \quad and \quad $e^{\ln x} = $ \ans{$x$}

\vspace{4pt}
If $f(x) = 3\ln x$, evaluate $f(e^4) = 3\ln(e^4) = 3 \cdot \ans{4} = $ \ans{12}

\vspace{4pt}
Solve $3\ln x = 9$: \quad $\ln x = $ \ans{3} \quad $\Rightarrow$ \quad $x = e^{\ans{3}} = $ \ans{$e^3$}
\end{notesbox}

\end{document}
